\chapter{Geltungsbereich und Grundbegriffe}

Dieser Leitfaden ergänzt Band~6 um eine ausgewählte Übersicht der
zentralen Definitionen, Charakterisierungen und Folgerungen für
\emph{injektive Funktionen zwischen Mengen}. Grundlage sind die
Mengenlehre in ZF und der allgemeine Funktionsbegriff aus Band~5.
Sämtliche Aussagen gelten ohne Auswahlaxiom. Beweise und
Beweisskizzen entfallen.

\section{Gegenstand und Abgrenzung}

Aufgenommen werden Aussagen, die Injektivität charakterisieren,
sie unter Konstruktionen erhalten oder aus ihr folgen.
Allgemeine Funktionsregeln werden vorausgesetzt; ihre
Definitionen erscheinen nur dort, wo sie für die Notation nötig sind.
Endlichkeit ist keine allgemeine Voraussetzung.

Surjektive und bijektive Funktionen, deren besondere Sätze und die
allgemeine Theorie der Umkehrfunktionen gehören in die Bände~7 und~8.
Ordnungs- und Äquivalenzrelationen, Mächtigkeitstheorie,
Auswahlprinzipien sowie algebraische, topologische und andere
zusätzliche Strukturen werden hier nicht behandelt.
Eine ausdrücklich genannte Mengenbedingung, etwa Disjunktheit
oder Nichtleere, ist eine Voraussetzung des jeweiligen Satzes.

Die Auswahl ist ein Leitfaden zu den wichtigsten elementaren
Resultaten, kein Anspruch auf eine vollständige Liste aller
bekannten oder ableitbaren Aussagen. Aufbau und Gestaltung folgen
den Leitfäden zu Band~3, Band~4, Band~5, Band~28 und Band~37.

\section{Injektivität und Notation}

\begin{definition}[Injektive Funktion]\label{b06guide:def:injection}
Eine Funktion $f\colon A\to B$ heißt \emph{injektiv}, wenn
\[
 \forall x,y\in A\,
       \bigl(f(x)=f(y)\Rightarrow x=y\bigr).
\]
Wir schreiben dafür $f\colon A\rightarrowtail B$.
Die Menge der injektiven Funktionen von $A$ nach $B$ ist
\[
 \operatorname{Inj}(A,B)
   =\{f\in B^A\mid f\text{ ist injektiv}\},
 \qquad B^A=\{f\mid f\colon A\to B\}.
\]
\end{definition}

Wie im Hauptband ist eine Funktion ihr Graph, also eine Menge
geordneter Paare. Eine zulässige Zielmenge gehört nicht als
zusätzlicher Bestandteil zum Graphen. Funktionswerte werden nur
für Elemente des angegebenen Definitionsbereichs verwendet.

\begin{definition}[Bild, Urbild und Faser als Hilfsnotation]
Für $f\colon A\to B$, $U\subseteq A$, $V\subseteq B$ und $b\in B$ sei
\[
 \begin{aligned}
 f[U]&=\{f(x)\mid x\in U\},\\
 f^{-1}[V]&=\{x\in A\mid f(x)\in V\},\\
 \operatorname{Fib}_f(b)&=f^{-1}[\{b\}].
 \end{aligned}
\]
Die Schreibweise $f^{-1}[V]$ bezeichnet eine Urbildmenge,
keine Umkehrfunktion. Ferner ist $\mathcal P(A)$ die Potenzmenge
von $A$, und
$U\mathbin{\triangle}W=(U\setminus W)\cup(W\setminus U)$
bezeichnet die symmetrische Differenz.
\end{definition}

\section{Wegweiser durch die Resultate}

Die folgende Übersicht unterscheidet Kriterien zum Erkennen
der Injektivität von Folgerungen und Konstruktionen.
Alle Voraussetzungen stehen in den jeweiligen Aussagen.

\begin{center}
\begin{tabularx}{\textwidth}{@{}p{.29\textwidth}X@{}}
\toprule
Fragestellung & Fundstelle\\
\midrule
Injektivität erkennen &
Punkt-, Faser- und Teilmengenkriterien in
Kapitel~\ref{b06guide:ch:criteria}.\\
Mengen zurückgewinnen &
Bild- und Urbildregeln in Kapitel~\ref{b06guide:ch:images}.\\
Kompositionen beurteilen &
Kürzen, Linksinversen und eindeutige Faktorisierung in
Kapitel~\ref{b06guide:ch:composition}.\\
Injektionen zusammensetzen &
Zusammenfügen und Erweiterungen in
Kapitel~\ref{b06guide:ch:extensions}.\\
Neue Injektionen bilden &
Produkte, Funktionenmengen und Mengenabbildungen in
Kapitel~\ref{b06guide:ch:constructions}.\\
\bottomrule
\end{tabularx}
\end{center}

Die Nummern mit Präfix $6H$ gehören ausschließlich zu diesem
Leitfaden. Die Nummerierung des vollständigen Fachbands bleibt
unverändert.